% Velocity from a position that only moves in whole counts.
\begin{tikzpicture}[font=\sffamily\scriptsize, x=1cm, y=1cm]

% ================================================================ counting method
\node[chlabel, anchor=west] at (0,3.7) {the counting method: how many counts arrived in one control period};

\foreach \r/\y/\lab/\n in {0/2.9/{fast: 50 counts per period}/{a velocity},
                           1/1.9/{slow: 1 count per period}/{a velocity, just},
                           2/0.9/{slower: 1 count every 4 periods}/{a quantiser}}{
  \node[signame] at (-0.3,\y) {\lab};
  \foreach \p in {0,1,2,3,4,5,6,7}{
    \draw[tick] (\p*0.9,\y-0.32) -- (\p*0.9,\y+0.32);
  }
}
% fast: counts in every period
\foreach \p in {0,1,2,3,4,5,6,7}{ \node[sample] at (\p*0.9+0.45,2.9) {}; }
\node[tlabel, anchor=west] at (7.3,2.9) {reported value is the motion};
% slow: one per period
\foreach \p in {0,1,2,3,4,5,6,7}{ \node[sample] at (\p*0.9+0.45,1.9) {}; }
\node[tlabel, anchor=west] at (7.3,1.9) {1 count per 1 ms: 15 rev/min at 4000 counts};
% slower: one every four periods
\foreach \p in {0,4}{ \node[sample] at (\p*0.9+0.45,0.9) {}; }
\node[tlabel, anchor=west] at (7.3,0.9) {reports 15, 0, 0, 0, 15, 0, 0, 0};
\draw[faultedge] (0.9,0.55) -- (0.9,0.9);
\draw[faultedge] (1.8,0.55) -- (1.8,0.9);
\draw[faultedge] (2.7,0.55) -- (2.7,0.9);
\node[tlabel, anchor=north, text=red!60!black] at (1.8,0.5) {three periods reporting zero while the shaft is turning};

% ================================================================ timing method
\node[chlabel, anchor=west] at (0,-0.6) {the timing method: the interval between edges, from chapter 3's capture};

\node[signame] at (-0.3,-1.4) {the same slow motion};
\draw[busline] (0,-1.4) -- (7.2,-1.4);
\foreach \x in {0.45,4.05}{ \node[sample] at (\x,-1.4) {}; }
\draw[tspan] (0.45,-1.9) -- node[tlabel]{4.000 ms, measured to the microsecond} (4.05,-1.9);
\node[tlabel, anchor=west] at (7.3,-1.4) {one count in 4.000 ms: 3.75 rev/min};

\node[signame] at (-0.3,-2.7) {and at high speed};
\draw[busline] (0,-2.7) -- (7.2,-2.7);
\foreach \x in {0.2,0.4,...,7.0}{ \node[sample] at (\x,-2.7) {}; }
\node[tlabel, anchor=west] at (7.3,-2.7) {edges arrive faster than they can be stamped};

% ================================================================ the crossover
\node[umlnote, text width=52mm, anchor=north west] at (0,-3.4)
  {Each method fails where the other works. The node uses the counting method
   above a threshold and the timing method below it, and the reported velocity
   carries which one produced it. The threshold is computed from the counts per
   revolution and the control period rather than chosen by taste: it is the
   speed at which one count per period is reached.};
\node[umlnote, text width=52mm, anchor=north west] at (5.6,-3.4)
  {The arithmetic is worth doing once here rather than discovering it in chapter
   16. With four thousand counts per revolution and a one millisecond period,
   one count per period is fifteen revolutions per minute. Below that the
   counting method reports fifteen or zero and nothing in between, and a
   controller differentiating that value will act on the quantisation.};
\end{tikzpicture}
